\documentclass{ltjsbook}
\usepackage{docmute} 
\usepackage{../myStyle} 

\begin{document} 
\chapter{統計環境Rの利用} 
\label{m04_usingR} 

\section{Rmarkdownによる文書作成}
\label{Rmarkdown}
前回はRに計算させたりデータを入力したり，と統計的な処理をいろいろさせてきました。RStudioは総合開発環境として，こうしたRの計算機能に加えて文書作成機能も持っています。ここではRStudioを使った文書作成を解説したいと思います。

表計算ソフトと文書作成ソフトは別々のものですがセットで販売・利用されることが多いものです。これらをまとめてオフィスソフトということがあります。Microsoftが提供するのはWord, Excel, PowerPointなどのOfficeシリーズですし，Appleが提供するのはPages, Numbers, Keynoteなどです。他にもフリーソフトウェアとしてLibreOfficeシリーズがあり，そこではWriter, Calc, Impress Presentationという名称になっていますが，何も文書作成，表計算，プレゼンテーションのアプリです。こうしたものが一般的に提供されているので，表計算ソフトで数値計算し，文書ファイルで報告書を書いて，プレゼンテーションアプリで紙芝居形式の発表をする，というのがよく行われています。これから皆さんもたくさんこうした資料を作ることになると思いますが，作成に当たっては，表計算ソフトで図表を作って，それをコピー＆ペーストして文書ファイルやプレゼンテーションファイルに切り貼りする，という方法を取ることでしょう。これが短い資料であれば問題ないのですが，長い資料や多くのデータになってきたりすると，貼り付ける図を間違えるとか，数字の修正をする前の図版だった(更新されたグラフにアップデートし忘れた)といった問題が生じやすくなります。これを\keytermL{コピペ汚染}{こぴぺおせん}と言ったりします。

文書作成とそこに含まれる図表や数字の計算が，同じ環境で行われればこの問題は無くなります。RStudioはそれを実現してくれる環境です。計算環境はRが持っていますから，それを文書に組み込む仕組みがあれば良いのです。RStudioではRmarkdownという書式を使ってこれに対応します。RmarkdownはMarkdownといわれるプレーンな文書作成形式にRを組み込んだRStudio独自のファイル形式です。メニューから\verb|File > New File > R Markdown|とすすむと，最終的にどの出力形式のファイルが必要か聞いてくるところがあります(図\ref{fig::05_08})。ここにあるように，最終的にHTML，PDF，Wordといったファイル形式で出力ができます\footnote{PDF形式に出力するには，TeXという別のシステムがまた必要ですが，これもRStudioからインストールすることができます。}。ここではひとまずHTMLを選びましょう。Titleには文書タイトルを，Authorには作者名を入れますので，ここも適当なタイトルと皆さんの名前を入れてOKを押してください。

\begin{figure}[ht]
    \centering
	\includegraphics[keepaspectratio,width=8cm]{../images/text05/Screen8.png}
    \caption{R Markdownからの最終出力を選択する}
    \label{fig::05_08}
\end{figure}

するとRmarkdown形式のファイルが開かれます。すでに何か書き込まれたファイルになっていますが，これは書き方の例が書いてあるのです。図\ref{fig::05_09}に最初の数行をキャプチャしたものを示しました。皆さんの手元の環境でも同じになっているか，確認しながら以下の解説を読んでください\footnote{もちろんTitleとAuthorはちがうはずですが！}。

\begin{figure}[ht]
    \centering
	\includegraphics[keepaspectratio,width=8cm]{../images/text05/Screen9.png}
    \caption{R Markdownからの最終出力を選択する}
    \label{fig::05_09}
\end{figure}

このファイルの1〜6行目は，文書全体の設定をしているところです。ここは基本的に変更する必要はありません。
1行目，6行目には`---`とありますが，この区切り記号が重要な意味を持ちますので，変更しないように特に注意してください。2行目はタイトル，3行目は作者，4行目は作成日，5行目は出力形式が書いてあります。これらは先程の図\ref{fig::05_08}で指定したものが反映されているはずです。これらは変更しても結構ですが，書式ルールから外れると後々うまくいかなくなるので注意してください。

7行目からが本文です。よく見ると，背景が灰色のゾーンと白色のゾーンがあると思います。灰色のゾーン(8〜10行目，18〜20行目など)は\keyterm{チャンク}{ちゃんく}{chunk}と呼ばれるRのコードを書く領域です。白色のゾーン(11〜17行目，21〜25行目など)は文書，本文を書くところになります。7行目以降はサンプルの文章が入っていますから，全て削除してしまっても構いません。実際に皆さんがレポート等を書く時は削除しましょう\footnote{後ほど指示しますが，本講の課題はRmarkdown形式のファイルを提出してもらいます}。でも今一旦そのままで，ファイルの上にある``Knit''というボタンを押してみましょう。これは\keyterm{ニット}{にっと}{knit}と言い，文章を編み上げるという意味です。このファイルが未保存の場合は，保存するファイル名を指定する画面が出ますが，任意の名前にして実行してみましょう\footnote{初めての場合，RStudioが\texttt{knit}に必要な関連パッケージをインターネットからダウンロードしようとします。ダウンロードしても良いか? と聞いてきますのでOKを選択してください。}

\begin{figure}[ht]
    \centering
	\includegraphics[keepaspectratio,width=12cm]{../images/text05/Screen10.png}
    \caption{R MarkdownファイルとKnitされたものの比較}
    \label{fig::05_10}
\end{figure}

図\ref{fig::05_10}には，Rmarkdownのファイルと出力結果を並べてみました\footnote{knitがうまく完了していれば，Viewerペインに表示されているはずです。}。対応関係がわかりますでしょうか。設定のところで入力したタイトルが大きく表示され，ついで作者名，日付が示されています。その後で本文が続きます。

出力の方では本文がR Markdownから始まっています。これはもとのファイルの方では\verb|##|で示されたところに対応していますね。Rmarkdownファイルでは，\verb|#|記号はコメントアウトではなく，見出しの階層を示しています。\verb|##|はレベル2，\verb|###|とするとレベル3と，\verb|#|の数で階層の深さを表現しているのです。階層が低いほど文字が大きく表示される仕組みです\footnote{レポートの書き方についてはこの授業で教える範囲を超えますので詳細は他の授業に譲りますが，アカデミックライティングの基本は構造化された文章であることは踏まえておいてください。科学論文や学術的な文書というものは，段落＞見出し＞小見出しといった階層構造を持っています。特に心理学論文の場合は，段落が「問題」「方法」「結果」「考察」の4つであることが決まっていて，独自のスタイルに勝手に変えることはできません。}。

見出しのついてない文章に関しては，そのまま本文として出力されています。注目すべきは，灰色の背景を持ったRの領域の箇所です。例えば元ファイルの19行目には，\texttt{summary(cars)}と書かれていますね。これはRのコマンドで，\texttt{cars}というオブジェクトの要約を，\texttt{summary}関数を使って表示せよ，というものになっています。\textbf{本文にはこの命令文しか書いてありませんが}，右の結果を見ると灰色の背景でどのような命令がなされたかを示すと同時に，結果が表示されています(speed変数とdist変数の記述統計量が表示されています)。また，元のファイルの27行目では\texttt{plot(pressure)}というコードが書かれています。これは\texttt{plot}関数によって図が描けという命令文なのですが，これもファイルの方に反映されているはずです。

このように，計算結果をコピー＆ペーストで文書に書き写すのではなく，明示するのは計算手順だけでよいのです。このようにすることで，この出力された文書を受け取る側は，どのような計算がなされ，どのような結果になるかを確認できます。そしてそこに\keytermL{コピペ汚染}{こぴぺおせん}は生じるはずがないのです！これは科学論文を作成するときに，方法を明示することと間違いがないことを担保しているという意味で，大変便利で画期的な方法であることがおわかりいただけたかと思います。

\paragraph{自分で書いてみる}
それでは元のファイルの7行目以下を全て削除し，自分なりにRmarkdownを書いてみましょう。
\texttt{code:\ref{code::05_08}}のように書くのを目標とします。
このとき，Rのコードを書くチャンク領域を挿入するには，図\ref{fig::05_11}に示した上のボタンを押して，Rを選択します。\texttt{code:\ref{code::05_08}}に書いてあるように，バッククォーテーションを3つ繋げて，中括弧でrと書き(\verb|```{r}|)，バッククォーテーションで閉じるという方法でも結構です。

チャンクには，設定や一時的な実行を行うためのボタンが表示されます(図\ref{fig::05_12})。実行ボタンを押すとknitではなくその下に実行結果が表示されるので，自分の意図した結果が示されているか，確認できます。

\begin{lstlisting}[language=c++,label={code::05_08},caption={Rmarkdownを書いてみる}]
---
title: "れんしゅう"
author: "小杉考司"
date: "2/5/2021"
output: html_document
---

# 計算をする

## 四則演算

### 足し算

足し算をします。

```{r}
2+3
```

### 引き算

引き算をします。

```{r}
5-2
```
\end{lstlisting}

\begin{figure}[ht]
    \centering
	\includegraphics[keepaspectratio,width=10cm]{../images/text05/Screen11.png}
    \caption{Rチャンクを挿入するボタン}
    \label{fig::05_11}
\end{figure}

\begin{figure}[ht]
    \centering
	\includegraphics[keepaspectratio,width=10cm]{../images/text05/Screen12.png}
    \caption{チャンクの設定や実行}
    \label{fig::05_12}
\end{figure}

うまくできましたでしょうか。コードや設定に間違いがあればknitは中断してしまいますが，エラーメッセージを読んで対応しましょう。RMarkdownについて，詳しくは\citet{高橋201805}が一冊を使って解説していますし，\citet{紀ノ定201806}や\citet{kosugi2019}なども参考になります\footnote{専修大学の国里ゼミと小杉ゼミでは，卒論をRmarkdownで書くことを推奨しています。再現性の観点から便利であることはもちろん，文書の規定に合わせたり引用文献リストを自動生成させたりする，専修大学専用パッケージ\texttt{senshuRmd}を共同開発しており，これを利用します。}。

\section{外部データの読み込み}
Rmarkdownを使うと，Rスクリプト(計算の指示書)とレポートの文章を1つにまとめ，knitすることで欲しい形式に出力できることがわかりました。ここで，スクリプトに描かれているのは計算の指示書だけであり，データをどのように加工するかの手続きが記載されています。言い換えると，データを変更しても計算手続きだけは決まった形で進めることができるのです。

もし定期的に同じようなデータが上がってきて，何度も同じようなレポートをしなければならない，というような仕事があったとします。これは仮定の話ではありますが，実際は社会に出て，毎月，毎年の集計をしなければならないようなことはよくあることです。この時いつも同じ計算をマウスでカチカチと繰り返しているのは，馬鹿らしいことですね。同じ手続き，同じ文章で数字だけ違うのであれば，データだけ変えて同じ指示書に流し込んでやればいいのです。Rmarkdownはまさにそういう時に向いています。

もちろんそれだけでなく，計算手続の記録をつけるという意味で科学的にも重要だという話でしたね。研究実践の例として，データが手に入ったらデータファイルは加工せず，加工の部分は全てRスクリプト/Rmarkdownでやることで計算結果の捏造が起きていないかどうかをチェックすることができるのです。研究で得られたデータは公開することで，いろいろな人にチェックしてもらおうというオープンサイエンスの動きも広がっています。

さて，このようにデータファイルが外部にある場合，それをRの中に取り込む必要があります。外部データの取り込みは，分析を始める前の手続きではありますが，初心者がトラブルに巻き込まれやすいシーンでもあります。手順を丁寧にみていきましょう。

% データが数件しかないのであれば，上で解説した方法で\texttt{data.frame}型に1つ1つ入力していけば良いのですが，Rのコードでそれを実行するとなると大変です。また実際の研究のシーンであれば，機械が大量にデータをファイルに出力するといっこともあるでしょう。Web調査を行った結果なども，表計算ソフトで利用できる形に出力されるのが一般的です。
% ここではそうしたファイルからデータを読み込む方法について説明します。

ところで，データファイルを自分で作ることもあるかもしれません。データを手入力してファイルを作る場合は，表計算ソフトを使う方が便利です。このときに注意すべきことは，\textbf{1つ載せるには1つのデータしか入れない}ということです。これからデータを色々加工していきますが，加工する前のデータはなるべくプレーンな状態であったほうが良いのです(すでに加工された食材を違う料理に利用するのは大変ですから！)。

図\ref{fig::05_13}には分析に不向きな，悪いデータファイル作成の例を示しました。これをまずはじっくり眺めて，どこがまずいのか考えてみてください。少なくとも7つの点で問題があります\footnote{残念ながら官公庁などが提供するデータの中には，こうしたよくない形式のファイルがよくみられ，データの活用ができにくくなっている実体があります。}。
\begin{figure}[ht]
    \centering
	\includegraphics[keepaspectratio,width=12cm]{../images/text05/Screen13.png}
    \caption{分析に不向きなデータファイル}
    \label{fig::05_13}
\end{figure}

分析しにくいよくないデータになっているのは次の箇所です。
\begin{enumerate}
	\item 1行目にこの表のタイトルのようなものが入っていますが，これは分析には使わない情報なので不要です。
	\item 1行目はセルを結合しています。データセットは\texttt{data.frame}型のように矩形であるべきなので，データのサイズを狂わせる行の存在は不適切です。
	\item 2行目には大学名が入っているようです。AからD列はC大学，FからIはD大学のデータのようですが，1つのデータセットに複数のデータが混在しているのはよくありません。
	\item 身長と体重，という変数名のところには1つのセルに複数の数字が入っています。カンマで区切られており，見た目にはわかりますが，機械的には(カンマのせいで)数字ではなく文字列として扱われます。データがデータになっていないのです。
	\item 変数2の列は「非常にそう思う」という言葉の回答と，数字が入っています。1つの列の中に複数の種類のデータが含まれるので分析できません(文字と数字の集計ができない。ちなみにこの場合は文字変数として扱われます)。
	\item 7行目に平均とありますが，データの平均はデータから計算できる量なので，Rで算出するものです。データの中に計算結果を入れておくべきではありません。
	\item I列6行目，D大学の杉田さんに関する変数3の箇所が空欄です。これはデータが入手できなかったところかもしれませんが，これは欠損値であることを明示すべきです。
	\item 「東京の大学」と「神奈川の大学」という複数のシートがあります。データファイルとしては1つのシートにすべての情報が入っていなければなりません。同時に分析できるデータなのであれば，一枚のシートにまとめるべきです。
\end{enumerate}
皆さんはいくつ気づいたでしょうか。これらの問題点を通じて言えることは「1つのデータセット，1つのセルには過不足ない情報を入れる」ということです。表を作って見せるのであればこのような入力でも良いのですが，ここでは数値データを分析することが目的なのですから，レイアウトや数字でない情報，計算してここから作れる情報は，分析に必要のない過剰な情報なのです。ちなみにアンダーラインや文字色，セルの色などは分析の時に自動的に受け落ちます(データではない情報なので)。

加えて分析の時に不便なのは，文字と数字の関係です。1つのセルに複数の数字が入る場合，区切り記号などをつけることが多いかと思いますが(スペースも機械にとっては文字です！)，こうした文字種・変数の種類・水準が混在しているのも情報過多ということになります。

今回の問題点を修正したのが図\ref{fig::05_14}のようなデータセットです。単一のデータセットで完結しています。また1つの列に一種類の変数，不要な情報はなく，過剰な情報もないデータセットになっているのがわかるかと思います。なお，変数3のところに\texttt{NA}とあるのは\texttt{No Answer}の略で\keyterm{欠損値}{けっそんち}{missing value}を意味しています。これは読み込む時に，欠損値をどのような文字にしているか，という指定ができますので問題になりません。

\begin{figure}[ht]
    \centering
	\includegraphics[keepaspectratio,width=12cm]{../images/text05/Screen14.png}
    \caption{分析に適したデータファイル}
    \label{fig::05_14}
\end{figure}

このように，データの入力の基本は，\textbf{1行に1ケースの情報が入っている，過不足のない1つのデータセットを作ること}である，と言えます。

このようにして作られたデータファイルは，UTF-8の文字コードによるcsv形式で保存するのが基本です($\to$第\ref{chapter2}章\ref{suffix}節)。ファイルを保存する時に，ファイル形式をしっかり指定し，拡張子も\texttt{.csv}にしましょう。csvファイルはASCIIファイルですし，文字コードも世界基準のUTF-8にしておけばどのPC空でもみることができます。ある一企業が開発した特定のOS，特定のアプリケーションを持っていなければデータを見ることができない，というのは科学的研究の観点からは不適切なのです。

\paragraph{データファイルを読み込む}
\label{readCSV}
さて，ではデータファイルがプロジェクトの中に保存できたとしましょう。このファイル名を\texttt{sample.csv}とします。このデータファイルを読み込むには，読み込むための関数を使います。Rが持っている基本関数である\texttt{read.csv}関数を使えば読み込めるのですが，ここではより便利なパッケージの関数を使ってみましょう。

ここではRmarkdownではなく，プレーンなRスクリプトファイルをつくって話を進めます。
\verb|FIle > New File > R Script|から新しくファイルを開くか，先ほどのRスクリプトファイルの下に\texttt{code:\ref{code::05_09}}のコードを書いて実行してください。

\begin{lstlisting}[language=c++,label={code::05_09},caption={CSVファイルを読み込む}]
library(tidyverse)
dat <- read_csv("sample.csv",na = "NA")
\end{lstlisting}

\begin{description}
    \item[1行目] \texttt{tidyverse}パッケージを読み込む。次の\verb|read_csv|関数がこのパッケージに含まれているからです。
	\item[2行目] \texttt{dat}オブジェクトに，\verb|read_csv|関数で読み込むファイルを代入します。関数の第一引数はファイル名，\verb|na|は欠損値の指定で，\texttt{NA}という文字列を欠損値として扱うことを宣言しています。
\end{description}

このようにすることで，外部ファイルを読み込むことができます。

% \begin{tcolorbox}[title=FAQ；ファイルの文字化け,fonttitle=\gtfamily\sffamily\bfseries,colframe=blue,colback=blue!3!]
% 　ここでWindowsユーザの多くがエラーに遭遇します。「資料として配布されたファイルが文字化けして読み込めません」という相談が毎年頻発します。

% 　理由はWindows特有の過保護な親切設計と，いまだに局所的な古い文字コードを使っていることにあります。
% 資料ファイルはcsv形式でダウンロードして使ってもらうようにしています。こちらは世界標準のUTF-8でエンコードしてあります。
% Windowsでこれをダウンロードすると，親切にも「ダウンロードしたら開くよね」といって勝手にExcelで開きます。みなさんもファイルの中身が見られたら安心ですね。でもそれを保存する時に，うっかり中身を触ってしまうようなことがあれば，文字コードを変えて上書き保存してしまいます。これで文字化けします。Excelなどで開かずに，直接Rで読み込むことがポイントです。

% 　また読込の時にエラーが生じるのであれば，WindowsのRにおける文字コードが勝手に日本ローカルに設定されて違う文字コードで読み込んでしまうことになります。
% そんな時は読み込み関数に\verb|locale=locale(encoding=”utf8")|をオプションとして追加すると良いでしょう。
% \end{tcolorbox}


\section{整然データ}




\section{グラフの文法}
RとRStudioを使って計算し，文書を作成するという基本的な使い方についての概説は以上です。それではこれから，より実践的な演習に入っていきます。

まずはデータをファイルから読み込むこと，そしてそれを描画するグラフの文法について紹介します。



\subsection{グラフの描画}
つづいてグラフの描画に進みましょう。第\ref{chapter3}章\ref{dataVisualization}節で触れたように，\textbf{データは図にする}のが基本中の基本ですし，この図をプログラミングすることが客観性，再現性の観点からも重要です。適切な図からは適切な情報の読み取りができますし，人為的でなくても不適切な図から間違った印象を受け取ることもありますから，なぜ・どういうつもりで図にするかを明文化できるようにしておきましょう。このことについての専門書，\citet{Kieran2018}の第1章は，知覚心理学や認知心理学の観点からグラフの書き方について論じられています。是非一度読んでみてください。

ここでは\keyterm{可視化の文法}{かしかのぶんぽう}{Grammar of Graphics}にしたがって作図をする，\texttt{ggplot2}パッケージ\citep{ggplot2}を使いながら説明します。
\paragraph{可視化の文法}
可視化の文法に沿った作図のコードは，次のように考えていきます。
\begin{enumerate}
    \item X軸Y軸にどのような変数を置くかを考える。
    \item X軸Y軸から出来上がるキャンバスに，点や線，何をプロットするのかを指定する。
    \item その上にモデルを加えて，元のデータとの対応を考えたりする。別のデータを重ねることもある。
\end{enumerate}
これを踏まえて，先程の\texttt{code:\ref{code::05_09}}の続きに\footnote{続きというのは，パッケージやデータの読み込みが終わっていることが必要だからです。描画に必要な\texttt{ggplot2}パッケージは，\texttt{tidyverse}パッケージを読み込むと同時に読み込まれています。}，\texttt{code:\ref{code::05_10}}を入力してください。

\begin{lstlisting}[language=c++,label={code::05_10},caption={まずはキャンパスを設定する}]
	g1 <- ggplot(data = dat,  mapping = aes(x = 身長, y = 体重))
	g1
\end{lstlisting}

\begin{description}
    \item[1行目] \texttt{g1}オブジェクトにggplotで作った図の情報を渡します。
	\item[2行目] \texttt{g1}って何? つまりこれを表示しろ，という命令です。\texttt{plot(g1)}と書いても結構です。
\end{description}
この1行目で，データは\texttt{dat}というオブジェクトを使い，x軸に\texttt{dat}オブジェクトの身長変数，y軸に体重変数を対応させたマッピングをしなさい，ということを指定しています。\texttt{aes}関数はピンとこないかもしれませんが，この1行目で，データは\texttt{dat}というオブジェクトを使い，x軸に\texttt{dat}オブジェクトの身長変数，y軸に体重変数を対応させたマッピングをしなさい，ということを指定しています。\texttt{aes}関数はなんと読むのかもピンとこない変な綴りだと思うかもしれませんが，実はaestheticの略です。エステティック，審美的・美学的思想という意味です。データの美的対応を考える箇所ですので，ちょうどいい名称ですね。この\texttt{aes}関数の中にはこのほかにも，色を何に対応させるか，サイズを何に対応させるか，群わけはどの変数に対応しているか，と言った対応関係を表す関数情報を設定するところだと思っておいてください。

さてこのように，データの縦横だけを指定したオブジェクトを作り，それを表示させると，プロットのところには灰色の背景を持ったキャンパスだけが表示されたと思います(図\ref{fig::05_15}の左上)\footnote{Macユーザであれば，x軸，y軸のラベル名が表示されず四角い豆腐のようなものしか出ないかもしれません。これは文字コードの問題で，次のコードを実行すれば直ります。\verb|old <- theme_set(theme_gray(base_family = "HiraKakuProN-W3"))|。あるいは，「Mac ggplot2 豆腐 文字化け」などで検索してみてください。}。

これがベースになって，次々と要素を付け加えていきます。次のコード\texttt{code:\ref{code::05_11}}，\texttt{code:\ref{code::05_12}}を書いて実行してみてください。

\begin{lstlisting}[language=c++,label={code::05_11},caption={キャンパスにデータ点をプロット}]
	g2 <- g1 + geom_point()
	g2
\end{lstlisting}

\begin{description}
    \item[1行目] \texttt{g1}オブジェクトに点(point)を打つ関数を追加したものを，\texttt{g2}オブジェクトとします。
	\item[2行目] \texttt{g2}って何? つまりこれを表示しろ，という命令です。\texttt{plot(g2)}と書いても結構です。
\end{description}

\begin{lstlisting}[language=c++,label={code::05_12},caption={線形モデルを当てはめる}]
	g3 <- g2 + geom_smooth(method="lm")
	g3
\end{lstlisting}

\begin{description}
    \item[1行目] \texttt{g2}オブジェクトに滑らかな線を引く関数を追加し，\texttt{g3}オブジェクトとします。線の関数は\texttt{lm}，Linear Modelの略で直線を引く時に使います。
	\item[2行目] \texttt{g3}って何? つまりこれを表示しろ，という命令です。\texttt{plot(g3)}と書いても結構です。
\end{description}


\begin{figure}[ht]
    \centering
	\includegraphics[keepaspectratio,width=16cm]{../images/text05/Rplot05_01.png}
    \caption{ggplotによる描画。左上は\texttt{code:\ref{code::05_10}}，右上は\texttt{code:\ref{code::05_11}}，左下は\texttt{code:\ref{code::05_12}}，右下は\texttt{code:\ref{code::05_13}の結果}}
    \label{fig::05_15}
\end{figure}


図がどのように変わっていくか，確認できたでしょうか。\texttt{g1}オブジェクトはただの白いキャンパスでしたが，それに点を打ったり(散布図)，線を引いたり(モデルの当てはめ)することでグラフがどんどん出来上がっていきます。
オブジェクトに次々追加していくことでこれができています。今回は理解のために逐一オブジェクトを作りましたが，慣れてくると一文で表現できます。

また今回使った\verb|geom_|という関数は，geometry(幾何学)の頭文字ですが，これを使ってさまざまなマッピングをします。よく使われるものに，次のようなものがあります。
\begin{itemize}
	\item \verb|geom_point| 点を打つ。散布図を書くことができる。
	\item \verb|geom_boxplot| ボックスプロットを書く。
	\item \verb|geom_bar| 棒グラフを書く。
	\item \verb|geom_line| 線を引く。折れ線グラフなどに使う。
	\item \verb|geom_smooth| 滑らかな線を引く。どのような線を引くかは色々な関数が考えられる。
\end{itemize}
これらの組み合わせで，さまざまなグラフが表現できるのです。少し書き換えるだけですぐにできますので，色々変えて遊んでみましょう。

最後に，次のコード\texttt{code:\ref{code::05_13}}を実行してみましょう。
\begin{lstlisting}[language=c++,label={code::05_13},caption={地区ごとに分割}]
	g4 <- g3 + facet_wrap(~地域)
	g4
\end{lstlisting}

\begin{description}
    \item[1行目] \texttt{g3}オブジェクトを地域変数で区切った図にせよ，という対応を指示。
	\item[2行目] \texttt{g4}を表示しろ，という命令です。 
\end{description}

今度は\verb|facet_wrap|という関数が出てきました。\verb|facet_|というのは分割，区切りを考える関数です。データにはさまざまな情報を書き込むことができますが，適当なグループで書き分けることでわかることも出てきます。今回は地域という変数で分割せよ，という表現をしており，その結果が図\ref{fig::05_15}の右下のようになっています。

可視化の文法については，習うより慣れろといった側面があります。さまざまな出力例をみながら，どのようにこのグラフが書けたのか，文法を読み解く練習もしてみましょう。ggplotによる出力例は，\url{https://www.r-graph-gallery.com/ggplot2-package.html}などが参考になります。

\section{課題}
今回の課題は，Rmarkdownファイルで作成してください。タイトルは「課題5」とし，作者明(Author)のところに学籍番号，氏名を記入してください。出力形式はHTMLを選んでおいていただければ結構です。\textbf{提出するファイルはエラーのない完備な.Rmdファイルです}。ファイル形式の間違い，出力にエラーが含まれている(未完成な)ものは返却しますので注意してください。
\paragraph{記述統計量を計算する}次の手順を実行し，データの記述統計量を計算してください。

\begin{enumerate}
	\item サンプルデータ\texttt{baseball2020.csv}をダウンロードし，プロジェクトフォルダに保存します。
	\item \texttt{code:\ref{code::05_exe1}}を実行し，データを整えます。
	\item データ全体の要約をする関数\verb|summary|を使って，データオブジェクトの記述統計量を表示してください。
	\item 身長変数heightの平均値を算出してください。平均値の関数は\verb|mean|です。
	\item 体重変数weightの標準偏差を算出してください。標準偏差の関数は\verb|sd|です。
	\item 年収変数salaryの中央値を算出してください。中央値の関数は\verb|median|です。
	\item 身長と体重の相関係数を算出してください。相関係数の関数は\verb|cor|で，2つの引数をとります。
\end{enumerate}

\begin{lstlisting}[language=c++,label={code::05_exe1},caption={データの読み込みと整え}]
	library(tidyverse)
	dat <- read_csv("baseball2020.csv")
	dat$team <- as.factor(dat$team)
	dat$position <- as.factor(dat$position)
	dat$bloodType <- as.factor(dat$bloodType)
\end{lstlisting}

\begin{description}
    \item[1行目] \texttt{tidyverse}パッケージを読み込みます。
	\item[2行目] データファイルを読み込みます。
	\item[3-5行目] \texttt{team,position,bloodType}などの変数は，名義尺度水準です。これが名義尺度水準であると設定し，同じ変数名に上書きしています。名義尺度水準にする関数は\texttt{as.factor}です。
\end{description}

\paragraph{図を書いてみる1}同じデータセットを使って，\texttt{code:\ref{code::05_exe1}}を実行してください。
\begin{lstlisting}[language=c++,label={code::05_exe2},caption={さまざまなグラフを書いてみる}]
	g1 <- ggplot(data = dat, mapping = aes(x = height, y = weight, color= team))
	g1 + geom_point()
	g1 + geom_smooth() 
	g2 <- ggplot(data = dat, mapping = aes(x = team, y = salary, fill= team))
	g2 + geom_boxplot()
\end{lstlisting}

\begin{description}
    \item[1行目] データは\texttt{dat}オブジェクト，x軸は身長変数，y軸は体重変数で，チームごとに色分けをします。
	\item[2行目] 散布図を書いてみます。
	\item[3行目] データに合うような線を引いてみます。
	\item[4行目] データは\texttt{dat}オブジェクト，x軸にチームを，y軸に年収変数をおいて，チームごとに色で塗りつぶします。
	\item[5行目] 箱ひげ図を書いてみます。
\end{description}
\paragraph{図を書いてみる2}図\ref{fig::05_16}と同じ図を書いてください。
\begin{figure}[ht]
    \centering
	\includegraphics[keepaspectratio,width=16cm]{../images/text05/Rplot05_02.png}
    \caption{課題の図}
    \label{fig::05_16}
\end{figure}

\end{document} 
